\documentclass[11pt,a4paper]{article}
\newcommand{\fichetype}{Exercices}
% Thème A4 commun, librement modifiable.
\usepackage[T1]{fontenc}
\usepackage[utf8]{inputenc}
\usepackage[french]{babel}
\usepackage{lmodern,amsmath,amssymb,mathtools,tikz,pgfplots}
\usepackage[margin=2.2cm,headheight=15pt]{geometry}
\usepackage{xcolor,fancyhdr,enumitem,hyperref,needspace}
\usetikzlibrary{arrows.meta,positioning,calc}
\pgfplotsset{compat=1.18,width=11cm,height=6cm}
\definecolor{Ink}{HTML}{202124}
\definecolor{Blue}{HTML}{356B9E}
\definecolor{Muted}{HTML}{737579}
\color{Ink}
\renewcommand{\familydefault}{\sfdefault}
\setlength{\parindent}{0pt}
\setlength{\parskip}{6pt}
\setlist{itemsep=3pt,topsep=4pt}
\pagestyle{fancy}\fancyhf{}
\fancyhead[L]{\small\color{Muted}MATHS LYCÉE}
\fancyhead[R]{\small\color{Muted}2026--2027}
\fancyfoot[L]{\small\color{Muted}\fichetype}
\fancyfoot[R]{\small\thepage}
\renewcommand{\headrulewidth}{0.3pt}
\hypersetup{colorlinks=true,urlcolor=Blue,linkcolor=Blue}
\newcommand{\R}{\mathbb{R}}
\newcommand{\N}{\mathbb{N}}
\newcommand{\E}{\mathbb{E}}
\newcommand{\V}{\mathrm{V}}
\newcommand{\exercice}[3]{\Needspace{6\baselineskip}\vspace{0.5cm}{\color{Blue}\small\textbf{EXERCICE #1}\hfill #2}\par{\large\bfseries #3}\par}

\begin{document}
{\small\color{Blue}PREMIÈRE SPÉCIALITÉ}\par{\LARGE\bfseries Simulation et échantillonnage}\par{\large Exercices}\par\vspace{0.2cm}\hrule\vspace{0.4cm}
Les exercices progressent de l’application directe vers la recherche. Détailler les calculs et justifier les réponses.
\par\noindent\begin{minipage}{\linewidth}\exercice{1}{Application}{Fréquence}
Une simulation donne $37$ succès en $80$ essais. Calculer la fréquence et expliquer pourquoi elle peut différer de $p=0,4$.
\end{minipage}\par
\par\noindent\begin{minipage}{\linewidth}\exercice{2}{Application}{Une loi simulée}
Un programme renvoie $1$ si $u<0,3$, et $0$ sinon, où $u$ est uniforme sur $[0,1[$. Donner la loi et l’espérance.
\end{minipage}\par
\par\noindent\begin{minipage}{\linewidth}\exercice{3}{Application}{Moyenne}
Pour l’échantillon $(2,6,3,1,4,5)$, calculer la moyenne et la fréquence de la valeur $6$.
\end{minipage}\par
\par\noindent\begin{minipage}{\linewidth}\exercice{4}{Entraînement}{Déboguer}
Un élève tire une fois un dé puis remplit une liste de cent copies de ce résultat. Pourquoi son échantillon est-il incorrect ?
\end{minipage}\par
\par\noindent\begin{minipage}{\linewidth}\exercice{5}{Entraînement}{Échelle}
Pour $\sigma=3$, calculer $\sigma/\sqrt n$ pour $n=100$ et $n=400$. Interpréter.
\end{minipage}\par
\par\noindent\begin{minipage}{\linewidth}\exercice{6}{Entraînement}{Écart normalisé}
Une variable a $\mu=10$, $\sigma=2$. Pour $n=64$, on observe $\overline x=10,5$. Calculer l’écart normalisé.
\end{minipage}\par
\par\noindent\begin{minipage}{\linewidth}\exercice{7}{Entraînement}{Seuils cumulés}
Proposer les seuils d’une simulation d’une variable prenant $-2$, $1$ et $4$ avec probabilités $0,25$, $0,35$, $0,40$. Calculer son espérance.
\end{minipage}\par
\par\noindent\begin{minipage}{\linewidth}\exercice{8}{Approfondissement}{Conclusion expérimentale}
Une expérience montre $946$ moyennes dans une bande donnée parmi $1000$ échantillons. Peut-on affirmer que la probabilité théorique est exactement $0,946$ ?
\end{minipage}\par
\end{document}
